\chapter{GIỚI THIỆU}
\label{chap:intro}

\section{Bài toán sinh mô tả ảnh}

Sinh mô tả ảnh (image captioning) là bài toán nhận một bức ảnh và sinh ra một
câu văn tự nhiên mô tả nội dung của ảnh. Khác với phân loại ảnh --- nơi đầu ra
là một nhãn rời rạc --- đầu ra ở đây là một \emph{chuỗi} từ có cấu trúc ngữ pháp,
đòi hỏi mô hình vừa ``nhìn hiểu'' ảnh (thị giác) vừa ``nói được'' (ngôn ngữ).
Bài toán vì vậy là ví dụ tiêu biểu của học đa phương thức (multimodal learning),
nơi hai phương thức thị giác và ngôn ngữ phải được kết nối trong cùng một mô hình
\autocite{vinyals2015show,xu2015show}.

\section{Vị trí đề tài trong sơ đồ kiến trúc đa phương thức của môn học}

Sơ đồ kiến trúc đa phương thức tổng quát mà giảng viên giới thiệu gồm: các bộ
mã hóa riêng cho từng phương thức (image encoder, text encoder), một khối
\emph{hợp nhất} (fusion) kết nối các phương thức, và phần đầu ra rẽ thành hai
nhánh --- \emph{phân loại} (classification head) hoặc \emph{sinh chuỗi}
(generation head).

Ở đồ án giữa kỳ, tôi đã đi nhánh phân loại với bài toán Medical VQA: ảnh và câu
hỏi được mã hóa, hợp nhất, rồi phân loại ra câu trả lời. Đồ án cuối kỳ này đi
nốt nhánh còn lại: khối fusion chính là \emph{cross-attention} giữa chuỗi từ đang
sinh và 49 vùng đặc trưng của ảnh, còn đầu ra là decoder tự hồi quy sinh từng từ.
Hai đồ án ghép lại phủ trọn sơ đồ: cùng một meta-kiến trúc, chỉ thay đầu ra.

\section{Mục tiêu và câu hỏi nghiên cứu}

Đồ án đặt bốn câu hỏi, tương ứng bốn trục thí nghiệm:
\begin{enumerate}
	\item Trên tập dữ liệu nhỏ (Flickr8k), decoder \textbf{LSTM+attention} hay
	      \textbf{Transformer} học tốt hơn? Đường cong học của chúng khác nhau ra sao?
	\item Cơ chế \textbf{attention} đóng góp bao nhiêu vào chất lượng mô tả
	      (so sánh LSTM có/không attention)?
	\item Chiến lược giải mã \textbf{greedy} so với \textbf{beam search} thay đổi
	      chất lượng thế nào?
	\item (Mở rộng) Tinh chỉnh bằng học tăng cường \textbf{SCST} sau huấn luyện
	      teacher forcing được gì và mất gì?
\end{enumerate}

Ràng buộc tự đặt: \emph{tự cài đặt} toàn bộ decoder (không dùng
\texttt{nn.TransformerDecoder} đóng hộp, không dùng mô hình captioning huấn luyện
sẵn như BLIP/GIT), để hiểu và trình bày được từng thành phần kiến trúc.

\section{Đóng góp và cấu trúc báo cáo}

Báo cáo gồm 5 chương. Chương~\ref{chap:background} trình bày cơ sở lý thuyết:
khung encoder--decoder, các dạng attention, teacher forcing và SCST.
Chương~\ref{chap:methods} mô tả dữ liệu Flickr8k, tiền xử lý, và kiến trúc hai
decoder cùng cấu hình huấn luyện. Chương~\ref{chap:experiments} trình bày kết quả
định lượng (BLEU/CIDEr), đường cong học, và phân tích định tính qua bản đồ nhiệt
attention. Chương~\ref{chap:conclusion} kết luận và nêu hạn chế.
